\section{动量涂抹：基本思想}
\label{sec:basic}
\begin{figure}
\includegraphics[width=0.48\textwidth]{momsmear_basic_principle_14pt.pdf}
\caption{以 $d=1$ 维空间中的高斯波函数为例，比较传统涂抹与动量涂抹。
动量 $k$ 使动量空间中分布的中心发生平移，
从而导致位置空间中的振荡行为。}
\label{fig:transform}
\end{figure}

如上所述，在强子源或汇中使用夸克涂抹对于格点模拟至关重要，
它可以增大与目标物理态的重叠，这反映出强子是延展客体，
而非点状粒子。涂抹算符 $F$ 在时间上是对角的，对自旋是平凡的，并
作用在夸克场的位置与颜色指标上：
\begin{equation}
	(F q)_{\vect{x}} = \sum_{\vect{y}\in (a \mathbb{Z})^d}\, f_{\vect{x}-\vect{y}}\, \GC_{\vect{x}\vect{y}}\, q_{\vect{y}} \label{formwupN}\,,
\end{equation}
其中 $f$ 是标量函数，$\GC$ 是规范协变输运子——在自由情形下它是颜色与位置空间的单位矩阵——
$d$ 是空间维数，通常 $d=3$。
注意场 $q_{\vect{x}}$ 在格点上通常关于 $\vect{x}$ 周期，而 $f_{\vect{x}-\vect{y}}$ 不必关于 $\vect{x}-\vect{y}$ 周期。
在自由情形下，卷积式~\eqref{formwupN} 在傅里叶空间成为乘积：
\begin{equation}
	\sum_{\vect{x}\in (a \mathbb{Z})^d} e^{i\vect{p}\cdot\vect{x}}\,(F q)_{\vect{x}} =
	\tilde{f}(\vect{p})\, \tilde{q}_{\vect{p}}\,.	\label{wupfour}
\end{equation}
对于高斯的特殊情形，
\begin{equation}
\label{eq:smgauss}
	 f_{\vect{x}-\vect{y}} = f_{\vect{0}} \exp\left(-\frac{|\vect{x}-\vect{y}|^2}{2{\sigma}^2}\right)\,,
\end{equation}
傅里叶变换后的涂抹核仍是高斯型：
\begin{equation}
	\tilde{f}(\vect{p}) \equiv  \sum_{\vect{z}\in (a \mathbb{Z})^d} e^{i\vect{p}\cdot\vect{z}} f_{\vect{z}} = \tilde{f}({\vect{0}}) \exp\left( - \frac{\sigma^2\vect{p}^2}{2} \right)  \, .	\label{kerfour}
\end{equation}
因此，涂抹后的夸克算符与静止夸克 $\vect{p}=\vect{0}$ 具有最大重叠。
按照上述高斯动量分布，非零速度被压制。
显然，对于携带显著空间动量的强子，这样的涂抹可能适得其反。

找出了问题所在之后，很容易修改涂抹使其对运动强子表现良好。
我们希望夸克的动量分布在有限动量 $\vect{k}$ 周围居中，因此需要在动量空间中平移涂抹核：
\begin{equation}
	\tilde{f}(\vect{p}) \mapsto  \tilde{f}(\vect{p} - \vect{k}) \label{kerfourshift}\,,
\end{equation}
如图~\ref{fig:transform} 所示。
这转换为位置空间中的替换
\begin{equation}
	f_{\vect{z}} \mapsto e^{i\vect{k}\cdot\vect{z}} f_{\vect{z}} \ .
\end{equation}
于是我们的修改后的涂抹算符 $F_{(\vect{k})}$（其中 $F_{(\vect{0})}=F$）
可以形式上表示为
\begin{align}
	(F_{(\vect{k})} q)_{\vect{x}} &= \sum_{\vect{y}\in (a \mathbb{Z})^d} F_{(\vect{k})\vect{x}\vect{y}} q_{\vect{y}}\nonumber\\&=
	\sum_{\vect{y}\in (a \mathbb{Z})^d} e^{-i\vect{k}\cdot(\vect{x}-\vect{y})}\, f_{\vect{x}-\vect{y}}\, \GC_{\vect{x}\vect{y}}\, q_{\vect{y}} \label{formnewsm}\ .
\end{align}
唯一的新要素是附加的
相位因子 $\exp\left[-i\vect{k}\cdot(\vect{x}-\vect{y})\right]$。
注意夸克动量平移
$\vect{k}$ 不必是格点动量，即它不限于离散值 $\vect{k}\in(2\pi/L)\mathbb{Z}^d$。涂抹核 $f$
与依赖规范的因子 $\GC$ 可以取自任何现有涂抹方法。对于迭代涂抹方法，
附加相位因子可以容易地融入基本涂抹步骤。下面我们以
Wuppertal 涂抹为例予以说明。
原则上还可能存在其他效应，
例如波函数的洛伦兹收缩。这
已在例如文献~\cite{Roberts:2012tp,DellaMorte:2012xc} 中被研究过，
我们也将讨论这一可能性。

新的修正涂抹算符 $F_{(\vect{k})}$（式~\eqref{formnewsm}）继承了它所基于的涂抹算符 $F_{(\vect{0})}$ 的重要性质。
若未修正的涂抹算符是自伴的，则修正后的涂抹算符 $F_{(\vect{k})}$ 也是自伴的，
因为从 $f_{\vect{x}-\vect{y}} = f_{\vect{y}-\vect{x}}^*$ 和 $\GC_{\vect{x}\vect{y}} = \GC^\dagger_{\vect{y}\vect{x}}$ 可得
$F_{(\vect{k})\vect{x}\vect{y}} = F^\dagger_{(\vect{k})\vect{y}\vect{x}}$。
底层的涂抹算符 $F_{(\vect{0})}$ 应按立方群 $O_h$ 的不可约表示（irrep）变换。
通常这是平凡的 $A_1$ 表示，
但涂抹算符也可用于注入角动量
或添加非平庸的胶子激发~\cite{Lacock:1994qx}。
显然，具有动量平移 $\vect{k}\neq\vect{0}$ 的 $F_{(\vect{k})}$
破坏立方对称性。然而，当它与选择动量为 $\vect{p}$ 的强子的动量投影联合使用时，
只要 $\vect{k}\parallel \vect{p}$，
该涂抹算符仍处于对应于
动量方向 $\vect{p}$ 的 $O_h$ 小群的
$A_1$ 表示之中。
